\section{Related Work}
\label{sec:related}

\subsection{Bimanual Manipulation}
Bimanual manipulation has a long history in both classical control and, more recently, machine learning. The literature review of Smith \emph{et al.}~\cite{smith2012dualarm} organizes the two arm literature along two axes: coordination structure (leader-follower, symmetric, or asymmetric) and coupling representation (kinematic constraints, hybrid position/force, or trajectory parameterization). The vast majority of these systems achieve coordination through a shared state vector or a centralized controller. On the learning side, ALOHA / ACT~\cite{zhao2023aloha} showed that a single policy with full access to both arms' states and a shared visual feed can imitate bimanual tasks from a small dataset. Both lines of work share the assumption that ours deliberately violates: that the two arms have access to a common state, controller, or model. We treat the two arms as fully independent processes connected only by what the camera can see.

\subsection{Cooperative Manipulation Without Communication}
Our closest ancestor is the work of Wang and Schwager~\cite{wang2016forceants} on the Force-Amplifying N-Robot Transport System (Force-ANTS). They demonstrate that a team of mobile robots can cooperatively transport a heavy object without communication, by having each robot independently estimate the object's velocity and apply forces in the same direction. Our system differs in three areas: the shared signal is visual rather than force-based; the object geometry is exploited explicitly (the tray's two markers estimate tilt); and the task progresses through discrete coordination events, rather than a single continuous transport. The underlying principle, however, is the same in that each agent computes something locally from a signal that captures the team's joint state.

\subsection{Visual Servoing and Marker-Based Perception}
The runtime control loop is a position-based visual servo (PBVS)~\cite{chaumette2006visual1}: at each tick we measure the pose error of the gripper relative to a tray marker in the Cartesian frame and command the joints to reduce that error. We rely on fiducial markers~\cite{garrido2014aruco} for pose extraction, using OpenCV's ArUco module configured with the AprilTag 16h5 library~\cite{olson2011apriltag} for their easy to make and large codes. At our 0.3-0.8 m visual operating range, the per-marker pose estimate from \texttt{cv2.solvePnP} with \texttt{SOLVEPNP\_IPPE\_SQUARE} produces 5-10 mm 3D accuracy, which is better than the coordination thresholds the algorithms of \cref{sec:coordination} require.

\subsection{Open Hardware and Software Stacks}
We build on the LeRobot framework~\cite{cadene2026lerobot} for SO-101 motor control, calibration, and bus management. LeRobot provides a uniform Python interface across follower-arm hardware, which lets us treat each arm thread identically up to the port and calibration files. The two SO-101 arms are inexpensive, open-source designs, which makes modifying the hardware and getting started with the software much more accessible.